% =====================================================================
%  courses/aritmetica/indice.tex — índice detallado: aritmética preuniversitaria
% ---------------------------------------------------------------------
%  5 Unidades · 29 Temas
% =====================================================================

{
\hypersetup{linkcolor=AzulPizarra}
\tableofcontents
}
\newpage

\thispagestyle{fancy}

\begin{center}
  {\Large\sffamily\bfseries\color{AzulPizarra} Índice Temático — Aritmética Preuniversitaria}\\[4pt]
  {\small\sffamily\color{GrisMedio} 5 Unidades · 29 Temas · Teoría, Ejercicios y Resueltos}
\end{center}

\vspace{0.6cm}

% --- UNIDAD I --------------------------------------------------------------
\begin{tcolorbox}[enhanced, colback=AzulPizarra, colframe=DoradoAcadémico,
    arc=3pt, boxrule=1pt, left=8pt, right=8pt, top=5pt, bottom=5pt,
    fontupper=\sffamily\bfseries\color{white}]
  UNIDAD I \hfill Fundamentos Numéricos \hfill Temas 1–6
\end{tcolorbox}

\begin{multicols}{2}
  \begin{enumerate}[label=\textcolor{AzulPizarra}{\textbf{\arabic*.}},
      leftmargin=2em, itemsep=1pt]
    \item Sistemas de numeración
    \item Conjuntos numéricos
    \item Operaciones fundamentales
    \item Divisibilidad
    \item Números primos
    \item MCD y MCM
  \end{enumerate}
\end{multicols}

\vspace{0.4cm}

% --- UNIDAD II -------------------------------------------------------------
\begin{tcolorbox}[enhanced, colback=AzulMedio, colframe=DoradoAcadémico,
    arc=3pt, boxrule=1pt, left=8pt, right=8pt, top=5pt, bottom=5pt,
    fontupper=\sffamily\bfseries\color{white}]
  UNIDAD II \hfill Fracciones y Decimales \hfill Temas 7–13
\end{tcolorbox}

\begin{multicols}{2}
  \begin{enumerate}[label=\textcolor{AzulMedio}{\textbf{\arabic*.}},
      leftmargin=2em, itemsep=1pt, start=7]
    \item Fracciones
    \item Operaciones con fracciones
    \item Números decimales
    \item Razones y proporciones
    \item Magnitudes proporcionales
    \item Regla de tres simple
    \item Regla de tres compuesta
  \end{enumerate}
\end{multicols}

\vspace{0.4cm}

% --- UNIDAD III ------------------------------------------------------------
\begin{tcolorbox}[enhanced, colback=AzulPizarra, colframe=DoradoAcadémico,
    arc=3pt, boxrule=1pt, left=8pt, right=8pt, top=5pt, bottom=5pt,
    fontupper=\sffamily\bfseries\color{white}]
  UNIDAD III \hfill Porcentajes e Interés \hfill Temas 14–18
\end{tcolorbox}

\begin{multicols}{2}
  \begin{enumerate}[label=\textcolor{AzulPizarra}{\textbf{\arabic*.}},
      leftmargin=2em, itemsep=1pt, start=14]
    \item Tanto por ciento
    \item Interés simple
    \item Interés compuesto
    \item Descuentos
    \item Reparto proporcional
  \end{enumerate}
\end{multicols}

\vspace{0.4cm}

% --- UNIDAD IV -------------------------------------------------------------
\begin{tcolorbox}[enhanced, colback=AzulMedio, colframe=DoradoAcadémico,
    arc=3pt, boxrule=1pt, left=8pt, right=8pt, top=5pt, bottom=5pt,
    fontupper=\sffamily\bfseries\color{white}]
  UNIDAD IV \hfill Aplicaciones Comerciales \hfill Temas 19–24
\end{tcolorbox}

\begin{multicols}{2}
  \begin{enumerate}[label=\textcolor{AzulMedio}{\textbf{\arabic*.}},
      leftmargin=2em, itemsep=1pt, start=19]
    \item Ganancia y pérdida
    \item Mezclas
    \item Aleaciones
    \item Promedios
    \item Movimiento uniforme
    \item Trabajo conjunto
  \end{enumerate}
\end{multicols}

\vspace{0.4cm}

% --- UNIDAD V --------------------------------------------------------------
\begin{tcolorbox}[enhanced, colback=AzulPizarra, colframe=DoradoAcadémico,
    arc=3pt, boxrule=1pt, left=8pt, right=8pt, top=5pt, bottom=5pt,
    fontupper=\sffamily\bfseries\color{white}]
  UNIDAD V \hfill Conteo y Probabilidad \hfill Temas 25–29
\end{tcolorbox}

\begin{multicols}{2}
  \begin{enumerate}[label=\textcolor{AzulPizarra}{\textbf{\arabic*.}},
      leftmargin=2em, itemsep=1pt, start=25]
    \item Análisis combinatorio
    \item Permutaciones
    \item Variaciones
    \item Combinaciones
    \item Probabilidad
  \end{enumerate}
\end{multicols}

\vspace{0.6cm}

\begin{cajatip}
  Este curso desarrolla el pensamiento numérico, base de todos los cursos de
  ciencias. Se recomienda dominar la Unidad~I antes de avanzar, pues la
  divisibilidad y los conjuntos numéricos son herramientas transversales en
  exámenes de admisión.
\end{cajatip}

\newpage